% Sample / manual test for pcalendar. Build: make test   (XeLaTeX, font Parisa)
\documentclass[12pt]{article}

\usepackage{pcalendar}
\usepackage{pgffor}
\usepackage{xepersian}

\settextfont{Parisa}

\newcount\jd \newcount\wd
% \row{label}{content}: one line of a two-column listing
\newcommand\row[2]{\par\noindent\lr{#1}\quad #2}

% The conversion macros take a literal Y-M-D (a macro such as \d would not be expanded),
% so dates are spelled out through these helpers instead of a \foreach list.
\newcommand\GtoJ[1]{\row{#1}{\pdate{#1}\quad\pisodate{#1}}}
\newcommand\JtoG[1]{\row{#1}{\gdate{#1}\quad\gisodate{#1}}}
% pgf counts Monday as 0; pcalendarweekdayname counts Saturday as 0
\newcommand\WEEKDAY[1]{%
  \pcalendardatetojulian{#1}\jd
  \pgfcalendarjuliantoweekday{\jd}{\wd}%
  \advance\wd by 2 \ifnum\wd>6 \advance\wd by -7 \fi
  \row{#1}{\pcalendarweekdayname{\wd}}}

\begin{document}

\section*{\lr{pcalendar sample}}

\subsection*{Today}
\row{\texttt{\textbackslash pdate}}{\pdate{\year-\month-\day}}
\row{\texttt{\textbackslash pisodate}}{\pisodate{\year-\month-\day}}

\subsection*{Names}
\row{months}{\foreach \m in {1,...,12}{\pcalendarmonthname{\m}\space}}
\row{gregorian months}{\foreach \m in {1,...,12}{\pcalendargregorianmonthname{\m}\space}}
\row{weekdays}{\foreach \w in {0,...,6}{\pcalendarweekdayname{\w}\space}}
\row{short weekdays}{\foreach \w in {0,...,6}{\pcalendarweekdayshortname{\w}\space}}

\subsection*{Gregorian to Jalali (\texttt{\textbackslash pdate}, \texttt{\textbackslash pisodate})}
\GtoJ{2015-03-21}
\GtoJ{2021-03-20}
\GtoJ{2024-03-20}
\GtoJ{2024-12-31}
\GtoJ{2025-03-20}

\subsection*{Jalali to Gregorian (\texttt{\textbackslash gdate}, \texttt{\textbackslash gisodate})}
\JtoG{1394-01-01}
\JtoG{1399-12-30}
\JtoG{1403-01-01}
\JtoG{1403-10-11}
\JtoG{1403-12-30}

\subsection*{First day of each year (\texttt{\textbackslash gisodate})}
\foreach \y in {1380,1385,...,1405}{%
  \row{\y}{\hbox{\gisodate{\y-01-01}}}%
}

\subsection*{Julian day numbers}
\pcalendardatetojulian{1394-01-01}\jd
\row{1394-01-01}{\number\jd}
\pcalendardatetojulian{1394-01-01+10}\jd
\row{1394-01-01 +10}{\number\jd}
\pcalendarjuliantodate\jd\y\m\d
\row{back to Jalali}{\number\y-\number\m-\number\d}
\pcalendardatetojulian{1403-12-29+1}\jd
\pcalendarjuliantodate\jd\y\m\d
\row{1403-12-29 +1 (leap year, expect 1403-12-30)}{\number\y-\number\m-\number\d}
\pcalendardatetojulian{1403-12-29+2}\jd
\pcalendarjuliantodate\jd\y\m\d
\row{1403-12-29 +2 (expect 1404-1-1)}{\number\y-\number\m-\number\d}

\subsection*{Weekday of a Jalali date}
\WEEKDAY{1394-01-01}
\WEEKDAY{1403-01-01}
\WEEKDAY{1404-01-01}

\subsection*{Animal year names (\texttt{\textbackslash pcalendaryearname})}
\foreach \y in {1380,...,1403}{\y~\pcalendaryearname{\y}\quad}

\end{document}
